\scanpage{12}%
8.6. La suite exacte (8.5.9) de systèmes locaux sur $U_d$ correspondant aux suites exactes (8.5.1) met en évidence que le\deleted{s} systêmes loca\deleted{ux}\handwritten{l} \unsure{corres}\inserted{pondant \deleted{aux}\handwritten{au}} \deleted{\unsure{des}} $(E(Y)\otimes_W k)_{ss}$ \deleted{et aux $E(Y,\underline O_Y)_{ss}$ ont}\inserted{\handwritten{a}}, en plus d'un sous-système \deleted{resp. d'un système quotient} \inserted{\handwritten{du premier}} de nature essentiellement triviale (car provenant du corps de base, donc constant du point de vue géométrique) \deleted{\unsure{une}}\handwritten{\unsure{un}} \unsure{composant} \deleted{commun $\mathrm{Coim}\,\varphi=\mathrm{Im}\,\varphi$} \inserted{\handwritten{$E(Y,\underline O_Y)_{ss}$}}, dont on s'attend qu'il donne lieu à une représentation tout à fait non triviale de $G$ (NB si on n'aime pas les groupes fondamentaux généraux et les systèmes locaux, on peut tout interpréter en termes de la suite exacte 8.5.1 correspondant au \uline{point générique} $s$ de l'espace des hypersurfaces de degré $d$ dans $\mathrm P^r_k$ , et remplacer $G$ par le groupe $\mathrm{Gal}(\overline{k(s)}/k(s))$, - mais on perd alors les frobenius ...). On aimerait pouvoir préciser ce point. Par exemple, \inserted{\handwritten{Si $n\geq2$, cette représentation est-elle \uline{irréductible}, ou même}} l'image du groupe \unsure{fondamental} \inserted{\handwritten{$G$}} dans $\mathrm{Aut}(V)$ ($V$ étant le \deleted{\unsure{corps}} vectoriel de dimension finie sur $\underline{\mathrm F}_p$ où s'effectue la représentation envisagée de $G$) est-elle égale à $\mathrm{Aut}(V)$ tout entier ? \deleted{Si oui,} \inserted{\handwritten{(Ce n'est pas \deleted{\unsure{toujours vrai}} lorsque $X=\mathrm P_k^{r=1}$, cas qu'on peut traiter trivialement.) Lorsque la condition \unsure{ci-dessus} est remplie,}} alors dès que \inserted{\handwritten{\unsure{l'on a}}} $\dim V\geq2$ (ce qui est sans doute \deleted{\unsure{vérifié}}\handwritten{vrai} dès que $d$ est assez grand, $\dim V$ tendant vers $\infty$ avec $d$, cf. 8.5.8), tout élément de $\underline{\mathrm F}_p$ est trace d'un élément de $G$ opérant sur $V$. De ceci, \inserted{de la ``formule de Woodshole'' \deleted{\handwritten{\unsure{pour les points fixes}}}} et du théorème de Čebotareff, on déduirait alors, \deleted{\unsure{k}}\handwritten{(}si le corps $k$ est fini\unsure{)}, qu'on peut, pour n'importe quel ouvert non vide $U$ de l'espace des hypersurfaces de degré $d$ dans $\mathrm P^r_k$ ($d$ étant choisi comme dessus) trouver un point \deleted{\unsure{de}} $s$ de $U$ à coéfficients dans une extension finie $k'$ de $k$, définissant une section $Y$ telle que $\mathrm{card}\,Y(k')$ soit congru mod $p$ à un élément de $\underline{\mathrm F}_p$ donné à l'avance. Au moment d'écrire ces lignes, et malgré l'assistance de gens compétents comme Katz, il ne semble pas qu'un résultat de cette nature soit prouvé\deleted{,}\handwritten{(} même si $X=\mathrm P^{2\handwritten{=r}}$ i.e. pour des courbes planes \inserted{\handwritten{$Y$}} de grand degré\handwritten{)}. \deleted{Ce résultat serait déjà bien intéressant, car il impliquerait déjà une réponse affirmative à la conjecture du n\textsuperscript{o}6 sur le niveau de la cohomologie évanescente (dont il est question aussi dans 8.5.8).}
% archive sheet number: 11

